\documentclass[10pt,a4paper]{amsart}
\usepackage[margin=19mm]{geometry}
\usepackage{amsmath,amssymb,mathtools}
\usepackage[scr=rsfs,cal=euler]{mathalfa}
\usepackage{xfrac}
\usepackage[unicode,hidelinks,bookmarks=false]{hyperref}
\newcommand{\ie}{i.e.\ }
\newcommand{\cf}{cf.\ }
\newcommand{\resp}{respectively\ }
\newcommand{\Subsection}{\subsection}
\providecommand{\nobreakdash}{\nobreak\hbox{-}}
\setlength{\emergencystretch}{2em}
\multlinegap0pt
\usepackage{csquotes}
\newtheorem{lemma}{Lemma}
\title{Extreme value theory for geometric Brownian motion and pricing of short maturity options}
\author{Ng Ze-An}
\date{Source: arXiv:2505.08036v2 (CC BY 4.0)}
\begin{document}
\maketitle

\noindent\emph{Source and licence.} Extreme value theory for geometric Brownian motion and pricing of short maturity options. Version arXiv:2505.08036v2 (CC BY 4.0), distributed by arXiv under the Creative Commons Attribution 4.0 International licence, http://creativecommons.org/licenses/by/4.0/, as recorded in the arXiv licence metadata for this exact version and shown on the abstract page https://arxiv.org/abs/2505.08036v2. Full licence notice: \texttt{licenses/arxiv-2505.08036-CC-BY-4.0.txt}.\par\smallskip

\noindent\textit{Editorial scope.} Counted unit: the article's lemma lemma (source0.tex lines 208--213). The article's complete original proof of it is delivered with it (source0.tex lines 214--407, one \texttt{proof} environment of 194 lines). No mathematics is written, repaired or re-derived: the statement and the proof are the article's own lines, in the article's own notation, and no retained line is substituted or re-flowed.  No further source-local context is needed: every cross-reference of every retained line resolves inside the two delivered spans.  Nothing was omitted from any retained span, so the package carries no omission record.  No retained line contains a citation, so the wrapper carries no bibliography and no bibliography entry.  No retained line contains a figure environment, an image-inclusion command or drawing code, so no image is delivered and the package declares no asset dependency.  Editorial formatting only, in addition: an AMS \texttt{amsart} preamble that restates the article's own declarations and macros used by the retained lines, namely the environments \texttt{lemma} on separate counters, together with the standard packages those lines need.  The retained environments print in this document as Lemma 1 in order of appearance, whereas the article prints the same items with the numbering of its own full document.  The complete native source of the article as distributed in the arXiv e-print for this exact version is bundled unchanged as \texttt{sources/178008/source0.tex}.
\begin{lemma}We have
$$
\mathbb{E}_{\mathbb{Q}_{M-f(T), T}}\left[\sup _{0 \leq t \leq T}\left|X_{t}-e^{\frac{t}{T} \sigma M}\right|\right]=O(\sqrt{T})
$$
as $T \rightarrow 0^{+}$, with the implied constant depending only on $f, M$.
\end{lemma}

\begin{proof}Since $X$ is a geometric Brownian motion, it admits the explicit solution
$$
X_{t}=\exp \left(C t+\sigma W_{t}\right)
$$
where for convenience we have written $C:=\mu-\frac{\sigma^{2}}{2}$. Write $W_{t}=\frac{t}{T} W_{T}+B_{t}$, where
$$
B_{t}:=W_{t}-\frac{t}{T} W_{T}
$$
is a standard Brownian bridge, independent of $W_{T}$. We then have
$$
X_{t}=\exp \left(C t-\sigma B_{t}+\frac{\sigma t}{T} W_{T}\right)
$$

Let $D$ be the event $\left\{W_{T} \geq M-f(T)\right\}$. We compute
$$
\begin{aligned}
& \mathbb{E}_{\mathbb{Q}_{M-f(T), T}}\left[\sup _{0 \leq t \leq T}\left|X_{t}-e^{\frac{t}{T} \sigma M}\right|\right] \\
& \leq \mathbb{E}_{\mathbb{Q}_{M-f(T), T}}\left[\sup _{0 \leq t \leq T}\left|\exp \left(C t-\sigma B_{t}+\frac{\sigma t}{T} W_{T}\right)-e^{(\sigma t / T) W_{T}}\right|\right] \\
& +\mathbb{E}_{\mathbb{Q}_{M-f(T), T}}\left[\sup _{0 \leq t \leq T}\left|e^{(\sigma t / T) W_{T}}-e^{\frac{t}{T} \sigma M}\right|\right]
\end{aligned}
$$

By monotonicity, the supremum in the last term occurs at $t=T$, and hence the last term is of order $O(\sqrt{T})$ by Lemma 1. For the first term, we claim that
$$
\begin{aligned}
& \mathbb{E}_{\mathbb{Q}_{M-f(T), T}} {\left[\sup _{0 \leq t \leq T}\left|\exp \left(C t-\sigma B_{t}+\frac{\sigma t}{T} W_{T}\right)-e^{(\sigma t / T) W_{T}}\right|\right] } \\
& \leq \mathbb{E}_{\mathbb{Q}_{M-f(T), T}}\left[\sup _{0 \leq t \leq T}\left|\exp \left(C t-\sigma B_{t}\right)-1\right| e^{\sigma W_{T}}\right] 
\end{aligned}
$$

Indeed, we have trivially
$$
\begin{aligned}
& \left.\sup _{0 \leq t \leq T}\left|\exp \left(C t-\sigma B_{t}+\frac{\sigma t}{T} W_{T}\right)-e^{(\sigma t / T) W_{T}}\right|\right] \\
& \left.\leq \sup _{0 \leq t \leq T} \sup _{0 \leq r \leq T}\left|\exp \left(C t-\sigma B_{t}+\frac{\sigma r}{T} W_{T}\right)-e^{(\sigma r / T) W_{T}}\right|\right] \\
& \left.=\sup _{0 \leq t \leq T} \sup _{0 \leq r \leq T} e^{(\sigma r / T) W_{T}}\left|\exp \left(C t-\sigma B_{t}\right)-1\right|\right]
\end{aligned}
$$

Since $\sigma>0$, and $W_{T}>0$ on the event $D$, we have that for all $t$, the inner supremum is attained at $r=T$, whence Equation 2 follows.

Next, using the independence of $B_{t}$ from $W_{T}$ (for all $t \in 0 \leq t \leq T$ ), denoting by $\mathcal{W}$ the $\sigma$-algebra generated by $W_{T}$, we have
$$
\begin{aligned}
& \mathbb{E}_{\mathbb{Q}_{M-f(T), T}}\left[\sup _{0 \leq t \leq T}\left|\exp \left(C t-\sigma B_{t}+\frac{\sigma t}{T} W_{T}\right)-e^{(\sigma t / T) W_{T}}\right|\right] \\
& \leq \mathbb{E}_{\mathbb{Q}_{M-f(T), T}}\left[\sup _{0 \leq t \leq T} e^{\sigma W_{T}}\left|\exp \left(C t-\sigma B_{t}\right)-1\right|\right] \\
& \leq \frac{\mathbb{E}\left[\mathbf{1}_{D} e^{\sigma W_{T}} \sup _{0 \leq t \leq T}\left|\exp \left(C t-\sigma B_{t}\right)-1\right|\right]}{\mathbb{P}(D)} \\
& =\frac{\mathbb{E}\left[\mathbb{E}\left[\mathbf{1}_{D} e^{\sigma W_{T}} \sup _{0 \leq t \leq T}\left|\exp \left(C t-\sigma B_{t}\right)-1\right| \mid \mathcal{W}\right]\right]}{\mathbb{P}(D)} \\
& =\frac{\mathbb{E}\left[\mathbf{1}_{D} e^{\sigma W_{T}} \mathbb{E}\left[\sup _{0 \leq t \leq T}\left|\exp \left(C t-\sigma B_{t}\right)-1\right| \mid \mathcal{W}\right]\right]}{\mathbb{P}(D)} \\
& =\frac{\mathbb{E}\left[\sup _{0 \leq t \leq T}\left|\exp \left(C t-\sigma B_{t}\right)-1\right|\right] \mathbb{E}\left[1_{D} e^{\left.\sigma W_{T}\right]}\right.}{\mathbb{P}(D)} \\
& =\mathbb{E}\left[\sup _{0 \leq t \leq T}\left|\exp \left(C t-\sigma B_{t}\right)-1\right|\right] \mathbb{E}_{\mathbb{Q}_{M-f(T), T}}\left[e^{\sigma W_{T}}\right]
\end{aligned}
$$

We now claim that $\mathbb{E}\left[\left|\sup _{0 \leq t \leq T} \exp \left(C t-\sigma B_{t}\right)-1\right|\right]$ is of order $O(\sqrt{T})$, while $\mathbb{E}_{\mathbb{Q}_{M-f(T), T}}\left[e^{\sigma W_{T}}\right]$ is of order $O(1)$, whence the result would follow. To see the first claim, note that we have
$$
\begin{aligned}
& \mathbb{E}\left[\left|\sup _{0 \leq t \leq T} \exp \left(C t-\sigma B_{t}\right)-1\right|\right] \\
& \leq \mathbb{E}\left[\mid \sup _{0 \leq t \leq T} \exp \left(C t-\sigma B_{t}\right)-\exp \left(-\sigma B_{t} \mid\right)\right]+\mathbb{E}\left[\left|\sup _{0 \leq t \leq T} \exp \left(-\sigma B_{t}\right)-1\right|\right] \\
& \leq\left(e^{C T}-1\right) \mathbb{E}\left[\mid \sup _{0 \leq t \leq T} \exp \left(-\sigma B_{t}\right)\right]+\mathbb{E}\left[\left|\sup _{0 \leq t \leq T} \exp \left(-\sigma B_{t}\right)-1\right|\right] .
\end{aligned}
$$

Since the former term tends to 0 as $T \rightarrow 0$, it will thus suffice to show that
$$
\mathbb{E}\left[\mid \sup _{0 \leq t \leq T} \exp \left(-\sigma B_{t}\right)-1\right]
$$
tends to 0 . We estimate

\begin{equation}
\mathbb{E}\left[\sup _{0 \leq t \leq T}\left|\exp \left(-\sigma B_{t}\right)-1\right|\right] \leq \mathbb{E}\left[\left|\exp \left(\sup _{0 \leq t \leq T}-\sigma B_{t}\right)-1\right|\right]+\mathbb{E}\left[\left|\exp \left(\inf _{0 \leq t \leq T}-\sigma B_{t}\right)-1\right|\right] .
\end{equation}


We show in turn that both terms in Equation (3) are of order $O(\sqrt{T})$. For the first term, we note that since
$$
B_{t}=W_{t}-\frac{t}{T} W_{T}
$$
we have
$$
0 \leq \sup _{t \in[0, T]}-B_{t} \leq M_{T}+\left|W_{T}\right|
$$
where
$$
M_{t}:=\sup _{t \in[0, T]}-W_{t} .
$$

So by the Cauchy-Schwartz inequality,
$$
\begin{aligned}
\mathbb{E}\left[\mid \sup _{0 \leq t \leq T} \exp \left(-\sigma B_{t}\right)-1\right] & =\mathbb{E}\left[\sup _{0 \leq t \leq T} \exp \left(-\sigma B_{t}\right)-1\right] \\
& \leq \sqrt{\mathbb{E}\left[\exp \left(2 \sigma M_{T}\right)\right]} \sqrt{\mathbb{E}\left[\exp \left(2 \sigma\left|W_{T}\right|\right)\right]}-1
\end{aligned}
$$

By the reflection principle, $M_{T}=\left|W_{T}\right|$ in law, so
$$
\mathbb{E}\left[\sup _{0 \leq t \leq T} \exp \left(-\sigma B_{t}\right)-1\right] \leq \mathbb{E}\left[\exp \left(2 \sigma\left|W_{T}\right|\right)\right]-1
$$

Letting $\Phi$ denote the CDF of a standard normal random variable, by standard formulae, we have
$$
\begin{aligned}
\mathbb{E}\left[\exp \left(2 \sigma\left|W_{T}\right|\right)\right] & =2 e^{2 T \sigma^{2}} \Phi(2 \sigma \sqrt{T}) \\
& =(1+O(\sqrt{T})) e^{2 T \sigma^{2}} \\
& =(1+O(\sqrt{T}))(1+O(T)) \\
& =1+O(\sqrt{T})
\end{aligned}
$$
whence
$$
\mathbb{E}\left[\left|\exp \left(\sup _{0 \leq t \leq T}-\sigma B_{t}\right)-1\right|\right]=O(\sqrt{T})
$$
as claimed.

Now we deal with the second term in Equation 3. Since $\inf _{0 \leq t \leq T}-B_{t} \leq 0$ almost surely, we have $\exp \left(\sigma \inf _{0 \leq t \leq T}-W_{t}\right) \leq 1$, so the second term is
$$
\mathbb{E}\left[1-\exp \left(\sigma \inf _{0 \leq t \leq T}-B_{t}\right)\right]
$$

Hence, it will suffice to show that
$$
\mathbb{E}\left[\exp \left(\sigma \inf _{0 \leq t \leq T}-B_{t}\right)\right] \rightarrow 1
$$

Again, since $B_{t}=W_{t}-\frac{t}{T} W_{T}$, we have
$$
m_{T}-\left|W_{T}\right| \leq \inf _{0 \leq t \leq T}-B_{t} \leq 0
$$
where $m_{T}:=\inf _{0 \leq t \leq T}-W_{t}$. So
$$
\begin{aligned}
E\left[\exp \left(\sigma \inf _{0 \leq t \leq T}-B_{t}\right)\right] & \geq E\left[\exp \left(\sigma\left(m_{T}-\left|W_{T}\right|\right)\right)\right] \\
& =\mathbb{E}\left[\frac{1}{\exp \left(\sigma\left(-m_{T}+\left|W_{T}\right|\right)\right)}\right] \\
& \geq \frac{1}{\mathbb{E}\left[\exp \left(\sigma\left(-m_{T}+\left|W_{T}\right|\right)\right)\right]}
\end{aligned}
$$
where in the last line we have applied Jensen's inequality. Applying the Cauchy Schwartz inequality, we have
$$
\frac{1}{\mathbb{E}\left[\exp \left(\sigma\left(-m_{T}+\left|W_{T}\right|\right)\right)\right]} \geq \frac{1}{\sqrt{\mathbb{E}\left[\exp \left(-2 \sigma m_{T}\right)\right]} \sqrt{\mathbb{E}\left[\exp \left(2 \sigma\left|W_{T}\right|\right)\right]}}
$$

By the reflection principle, $-m_{T}=\left|W_{T}\right|$ in distribution, so
$$
\frac{1}{\sqrt{\mathbb{E}\left[\exp \left(-2 \sigma m_{T}\right)\right]} \sqrt{\mathbb{E}\left[\exp \left(2 \sigma\left|W_{T}\right|\right)\right]}} \geq \frac{1}{\mathbb{E}\left[\exp \left(2 \sigma\left|W_{T}\right|\right)\right]}
$$

Consequently, we have
$$
\mathbb{E}\left[1-\exp \left(\sigma \inf _{0 \leq t \leq T}-B_{t}\right)\right] \leq 1-\frac{1}{\mathbb{E}\left[\exp \left(2 \sigma\left|W_{T}\right|\right)\right]}
$$

Since
$$
\mathbb{E}\left[\exp \left(2 \sigma\left|W_{T}\right|\right)\right]=1+O(\sqrt{T})
$$
as proven earlier, we deduce
$$
\mathbb{E}\left[1-\exp \left(\sigma \inf _{0 \leq t \leq T}-B_{t}\right)\right] \leq 1-\frac{1}{1+O(\sqrt{T})}=O(\sqrt{T})
$$
as claimed.

On the other hand, the second claim follows from a stopping time argument and standard estimates. Indeed, write
$$
\tau_{T}:=\inf \left\{t>0 \mid W_{T} \geq M-f(T)\right\}
$$

We have
$$
\begin{aligned}
\mathbb{E}_{\mathbb{Q}_{M-f(T), T}}\left[e^{\sigma W_{T}}\right] & =\frac{\mathbb{E}\left[1_{D} e^{\sigma W_{T}}\right]}{\mathbb{P}(D)} \\
& =\frac{\mathbb{E}\left[\mathbb{E}\left[\left[1_{D} e^{\sigma W_{T}} \mid \mathcal{F}_{\tau}\right]\right]\right.}{\mathbb{P}(D)} \\
& =\frac{\mathbb{E}\left[1_{D} \mathbb{E}\left[\left[e^{\sigma W_{T}} \mid \mathcal{F}_{\tau}\right]\right]\right.}{\mathbb{P}(D)} .
\end{aligned}
$$

On $D$, we have $\tau_{T} \leq T$ almost surely. Thus by the strong Markov property, conditional on $\mathcal{F}_{\tau}, R_{t}:=W_{t+\tau}$ is a Brownian motion with initial value $R_{0}=W_{\tau}=M-f(T)$. Thus
$$
\begin{aligned}
\mathbb{E}\left[\left[e^{\sigma W_{T}} \mid \mathcal{F}_{\tau}\right]\right. & =\mathbb{E}\left[\left[e^{\sigma R_{t-\tau}} \mid \mathcal{F}_{\tau}\right]\right] \\
& =\left.\mathbb{E}\left[e^{\sigma R_{t-r}}\right]\right|_{r=\tau}
\end{aligned}
$$
where in the last equality we have applied the freezing lemma. We recognise $e^{\sigma R_{t-r}}$ as a log normal random variable with mean $\exp \left(M-f(T)+\frac{t-r}{2}\right) \leq \exp \left(M+|f(T)|+\frac{T}{2}\right)<\exp (M+1):=C$ for all small enough $T$, uniformly over all $0 \leq r \leq t$. Thus
$$
\begin{aligned}
\mathbb{E}_{\mathbb{Q}_{M-f(T), T}}\left[e^{\sigma W_{T}}\right] & =\frac{\mathbb{E}\left[\left.1_{D} \mathbb{E}\left[e^{\sigma R_{t-r}}\right]\right|_{r=\tau}\right]}{\mathbb{P}(D)} \\
& \leq C \frac{\mathbb{E}\left[1_{D}\right]}{\mathbb{P}(D)} \\
& =C
\end{aligned}
$$

Thus $\mathbb{E}_{\mathbb{Q}_{M-f(T), T}}\left[e^{\sigma W_{T}}\right]$ is of order $O(1)$ as claimed, and this concludes the proof.
\end{proof}

\end{document}
